\begin{description}
\item[(a)] In this scenario, the asteroid would fall directly towards the
  star. However, in order to simplify the calculations, it is possible to
  consider that the asteroid would be in a degenerate elliptical orbit. In
  that case, the semi-minor axis would be infinitesimally small and the
  asteroid would still practically be moving on a straight line. It is also
  important to notice that the focii would virtually be at the periapsis and % sic: "focii", "apapsis"
  apapsis points. The degenerate elliptical orbit is shown in the following
  figure: \hfill(2.0 points)

% FIGURE NEEDED: degenerate ellipse diagram (2024 theory solutions, page 5 of
% 40): flattened ellipse with the Star at the left focus, the Asteroid at the
% bottom, semi-major axes labelled a, the flattening "b -> 0" marked, and the
% swept regions coloured Area I (green), Area II (red), Area III (blue).

  The semi-major axis of this elliptical orbit can be obtained through the
  following expression, where $m$ is the mass of the asteroid:
  \begin{align*}
    E = \frac{mv^2}{2} - \frac{GMm}{d} &= -\frac{GMm}{2a} \\
    \frac{GMm}{2d} - \frac{GMm}{d} &= -\frac{GMm}{2a} \\
    \frac{1}{2d} &= \frac{1}{2a} \\
    \therefore\ d &= a
  \end{align*}
  \hfill(2.0 points)

  Now, it is possible to use Kepler's third law to find an expression for the
  period of the orbit:
  \[
    \frac{T^2}{d^3} = \frac{4\pi^2}{GM}
      \longrightarrow T = 2\pi d\sqrt{\frac{d}{GM}}
  \]
  \hfill(1.0 points)

  If the asteroid is initially moving towards the star, it sweeps area I to
  reach the star. Using Kepler's second law, it is possible to calculate how
  long it takes for that to happen:
  \[
    \frac{\Delta t}{T} = \frac{A_I}{A_{Total}}
      = \frac{\frac{\pi ab}{4} - \frac{ab}{2}}{\pi ab}
  \]
  \hfill(2.0 points)
  \[
    \Delta t = \left(\frac{1}{4} - \frac{1}{2\pi}\right)T
  \]
  \[
    \boxed{\Delta t = \left(\frac{\pi}{2} - 1\right)d\sqrt{\frac{d}{GM}}}
  \]
  \hfill(1.0 points)

  Note that the radius of the star is negligible in these calculations since
  it is significantly smaller than $d$.

\item[(b)] If the asteroid is initially moving away from the star, it sweeps
  areas II and III before reaching the star. It is possible to again use
  Kepler's second law to find the time interval:
  \[
    \frac{\Delta t}{T} = \frac{A_{II} + A_{III}}{A_{Total}}
      = \frac{\frac{ab}{2} + \frac{\pi ab}{4} + \frac{\pi ab}{2}}{\pi ab}
  \]
  \hfill(1.5 points)
  \[
    \Delta t = \left(\frac{3}{4} + \frac{1}{2\pi}\right)T
  \]
  \[
    \boxed{\Delta t = \left(\frac{3\pi}{2} + 1\right)d\sqrt{\frac{d}{GM}}}
  \]
  \hfill(0.5 points)
\end{description}
